\documentclass[11pt]{article}
\usepackage[UTF8]{ctex}
\usepackage{geometry}
\usepackage{booktabs}
\usepackage{array}
\usepackage{hyperref}
\geometry{margin=1in}

\title{DevGuard：小鼠胚胎扰动细胞的阶段特异发育正常性判定框架}
\author{DevGuard working draft}
\date{2026-05-05}

\begin{document}
\maketitle

\section*{拟定英文标题}

\textbf{DevGuard: stage-specific developmental normality boundaries for perturbed mouse embryonic cells}

\section*{一句话主张}

DevGuard 学习小鼠每个发育时间点和谱系的正常范围，并用 conformal 校准判断扰动细胞是仍然正常、发育延迟、发育提前、命运偏航，还是成为正常图谱外异常状态。

\section*{路线边界}

新版 DevGuard 是独立的小鼠发育扰动正常性判定项目。它不使用 RDEG pipeline、RDEG-derived node graph、OT transition、GRN dynamics 或 matrix-writeback outputs。旧版 DevVCell response-recovery 原型已经完整归档到：

\begin{verbatim}
Old/legacy_devvcell_response_recovery_20260505/
\end{verbatim}

当前工作稿中的 quick-mode 结果只用于软件和 schema 验证，不作为生物学结论。真实论文结果需要在 E-MTAB-6967、GSE212050、GSE123187 或其他独立公开小鼠扰动数据接入后重新运行。

\section*{核心问题}

给定一个小鼠发育细胞和一个扰动，DevGuard 判断该细胞是否仍处于对应发育阶段和谱系的正常范围内。若不在正常范围内，进一步判定其最符合以下哪一种状态：

\begin{enumerate}
  \item \texttt{within\_stage\_normal}：当前阶段同谱系正常；
  \item \texttt{developmental\_delay}：当前阶段异常，但更像早期同谱系；
  \item \texttt{developmental\_acceleration}：当前阶段异常，但更像晚期同谱系；
  \item \texttt{fate\_deviation}：当前阶段异常，但更像其他谱系；
  \item \texttt{abnormal\_off\_normal}：所有正常参考集均不通过。
\end{enumerate}

\section*{方法概述}

DevGuard 在 control mouse embryo 或 gastruloid time-course 中，按时间点和谱系建立正常参考集：

\[
N(t,l) = \{x_i: x_i \text{ is a control cell at time } t \text{ and lineage } l\}.
\]

每个正常参考集被划分为 training、calibration 和 heldout test 三部分。Training cells 用于学习 embedding 和 reference distribution；calibration cells 用于估计 nonconformity score 的 conformal p-value；heldout test cells 用于检查 control false-positive rate。

\subsection*{Embedding}

MVP 使用可审计的浅层 embedding：

\begin{enumerate}
  \item library-size normalization；
  \item \texttt{log1p}；
  \item high-variance gene selection；
  \item \texttt{TruncatedSVD} latent embedding。
\end{enumerate}

\subsection*{Nonconformity score}

当前实现两类 score。

\paragraph{KNN distance}
\[
s_{\mathrm{knn}}(x;t,l)=\frac{1}{k}\sum_{j\in KNN_k(x,N_{\mathrm{train}}(t,l))}\lVert z_x-z_j\rVert_2.
\]

\paragraph{Mahalanobis distance}
\[
s_{\mathrm{mah}}(x;t,l)=(z_x-\mu_{t,l})^\top \Sigma_{t,l}^{-1}(z_x-\mu_{t,l}).
\]

\subsection*{Conformal p-value}

对 query cell \(x\)，其相对于 stage-lineage reference \(N(t,l)\) 的 conformal p-value 为：

\[
p(x;t,l)=
\frac{1+\#\{s_i^{\mathrm{cal}}\geq s(x;t,l)\}}
{1+n_{\mathrm{cal}}}.
\]

若 \(p(x;t,l)\geq \alpha\)，则认为 \(x\) 仍在该 stage-lineage 正常范围内。当前 MVP 默认 \(\alpha=0.05\)。

\section*{Quick-mode 软件验证结果}

运行命令：

\begin{verbatim}
python scripts/devguard/run_devguard_pipeline.py --mode quick
\end{verbatim}

该命令生成 deterministic mouse gastruloid software fixture，并完整执行 prepare、reference building、conformal calibration、perturbed-cell classification、DTI 和 Figure 1--3 草图生成。

\subsection*{Quick fixture 规模}

\begin{center}
\begin{tabular}{ll}
\toprule
项目 & 数值 \\
\midrule
species & \textit{Mus musculus} \\
system & gastruloid \\
control time points & E6.5, E7.5, E8.5 \\
lineages & epiblast, mesoderm, neural \\
control cells & 1,350 \\
perturbed cells & 120 \\
genes & 90 \\
reference groups & 9 \\
\bottomrule
\end{tabular}
\end{center}

\subsection*{Heldout control false-positive calibration}

在 quick-mode fixture 中，normality reference quality table 共包含 18 条记录，即 9 个 stage-lineage groups 乘以 2 种 score methods。最大 heldout control false-positive rate 为：

\[
\max(\mathrm{FPR}_{\mathrm{heldout}})=0.0333.
\]

该数值低于默认 \(\alpha=0.05\)，说明 quick fixture 下 conformal boundary 的软件行为符合预期。该结论只说明 pipeline 和 calibration 逻辑可运行，不等价于真实公开数据上的统计结论。

\subsection*{Perturbed-cell 五类分类}

Quick-mode 中设置了一个示例扰动 \texttt{FGF8\_KO}，所有 perturbed cells 的观测时间点和谱系为 E7.5 mesoderm。为了验证五类分类逻辑，fixture 中分别构造了 5 组表达模式，每组 24 个细胞。

\begin{center}
\begin{tabular}{lll}
\toprule
fixture pattern & DevGuard class & cell count \\
\midrule
within\_stage\_normal\_like & within\_stage\_normal & 24 \\
delay\_like & developmental\_delay & 24 \\
acceleration\_like & developmental\_acceleration & 24 \\
fate\_deviation\_like & fate\_deviation & 24 \\
abnormal\_like & abnormal\_off\_normal & 24 \\
\bottomrule
\end{tabular}
\end{center}

Condition-level summary 因此得到五类各占 20\%：

\begin{center}
\begin{tabular}{lc}
\toprule
normality class & fraction \\
\midrule
within\_stage\_normal & 0.20 \\
developmental\_delay & 0.20 \\
developmental\_acceleration & 0.20 \\
fate\_deviation & 0.20 \\
abnormal\_off\_normal & 0.20 \\
\bottomrule
\end{tabular}
\end{center}

\subsection*{Developmental Tolerance Index}

Developmental Tolerance Index 定义为：

\[
\mathrm{DTI}
=R_{\mathrm{normal}}
-R_{\mathrm{delay}}
-R_{\mathrm{accel}}
-R_{\mathrm{deviation}}
-R_{\mathrm{abnormal}}.
\]

Quick-mode 的 E7.5 mesoderm \texttt{FGF8\_KO} 示例中，每一类比例均为 0.20，因此：

\[
\mathrm{DTI}=0.20-0.20-0.20-0.20-0.20=-0.60.
\]

这个值用于验证 DTI 汇总代码是否按定义运行。真实论文中 DTI 应在真实 perturbation conditions、time points 和 lineages 上解释。

\section*{当前代码交付}

当前仓库已经包含以下新版 DevGuard 组件：

\begin{verbatim}
config/devguard/
scripts/devguard/
src/devguard/
docs/DEVGUARD_*.md
manuscript/main_cn.tex
tests/test_conformal.py
tests/test_classification.py
tests/test_tolerance.py
tests/test_io_contract.py
\end{verbatim}

关键输出文件由 quick mode 生成到：

\begin{verbatim}
data/processed/devguard/devguard_quick_mouse.h5ad
results/devguard/normality_reference/
results/devguard/perturbation_classification/
results/devguard/tolerance_index/
results/devguard/figures/
\end{verbatim}

其中 \texttt{data/} 和 \texttt{results/} 按仓库规则仍被 git ignore，不提交大数据和生成结果。

\section*{Figure 草图}

Quick mode 已生成以下草图：

\begin{enumerate}
  \item Figure 1：DevGuard conformal developmental normality 方法示意图；
  \item Figure 2：perturbed-cell 五类分类比例；
  \item Figure 3：Developmental Tolerance Index heatmap。
\end{enumerate}

这些文件位于：

\begin{verbatim}
results/devguard/figures/
\end{verbatim}

\section*{测试结果}

当前单元测试命令：

\begin{verbatim}
python -m pytest
\end{verbatim}

测试结果：

\begin{verbatim}
9 passed
\end{verbatim}

覆盖范围包括 conformal p-value、classification priority、DTI formula、quick fixture obs schema 和 RDEG legacy input rejection。

\section*{下一步真实数据接入}

Quick-mode 之后的真实论文工作应按以下顺序推进：

\begin{enumerate}
  \item 接入 E-MTAB-6967 / MouseGastrulationData，作为 main normality reference；
  \item 接入 GSE212050，使用 individual organoid barcode 做 leave-one-organoid-out false-positive control；
  \item 接入 GSE123187，做 spatial/tomo axis validation；
  \item 若小鼠 perturbation 数据不足，再考虑补充 human GSE208369 作为 supplementary cross-species demonstration，但不进入主图核心证据链；
  \item 每个真实 perturbation 结论均需与 heldout control false-positive rate 和 organoid heterogeneity baseline 对比。
\end{enumerate}

\section*{稿件写作边界}

正式文章中应使用以下关键词：

\begin{itemize}
  \item independent public mouse perturbation datasets；
  \item stage-specific developmental normality；
  \item conformal calibration；
  \item perturbed-cell classification；
  \item developmental tolerance index；
  \item organoid heterogeneity control；
  \item spatial/tomo-seq validation。
\end{itemize}

正式文章中不应写“基于 RDEG”、“继承 RDEG 的 OT”、“用 GRN 动力学推演扰动”、“从第一篇图谱流形出发”或“虚拟细胞生成未来表达矩阵”。

\end{document}
